\documentclass[12pt,a4paper]{article}

% ================================================================
% PAGE SETUP 
% ================================================================

\usepackage[a4paper,left=2cm,right=2cm,top=2cm,bottom=2cm]{geometry}
\usepackage{setspace}
\usepackage{changepage}
\usepackage{titlesec}

% ================================================================
% TYPESETTING AND LANGUAGE
% ================================================================

\usepackage[T1]{fontenc}
\usepackage{newtxtext,newtxmath}
\usepackage[ngerman]{babel}
\usepackage{csquotes}
\usepackage[normalem]{ulem}
\usepackage[colorlinks=true,linkcolor=blue,urlcolor=blue,citecolor=blue]{hyperref}
\newcommand{\ulhref}[2]{\href{#1}{\uline{#2}}}

% ================================================================
% BODY
% ================================================================

\begin{document}

\begin{center}
    \Large{
        \textbf{
            \uline{D178 PRAKTIKUM}
        }
    }
\end{center}

\begin{flushright}
    Quin Barnard \textit{27902277}\\
    Bianca van der Walt \textit{28512979}
\end{flushright}

\large{\textbf{Allgemein}}
\begin{itemize}
    \item Brauchen Lautsprecher und Beamer
    \item Gedrückte Handouts
    \item GA: Die TN werden in dreiköpfigen Gruppen arbeiten; überlegen wie sie sich im Raum 545 arrangieren werden
    \item Die Praktika könnten entweder von \textit{„Die Bremer Stadtmusikanten“} oder von \textit{„Rotkäppchen“} handeln. Die KL können sich entscheiden
    \item \ulhref{https://www.grimmstories.com/de/grimm_maerchen/die_bremer_stadtmusikanten}{Die Bremer Stadtmusikanten}
    \item \ulhref{https://www.grimmstories.com/de/grimm_maerchen/rotkaeppchen}{Rotkäppchen}
\end{itemize}

\large{\textbf{Inhalt und Ziel}}
\begin{itemize}
    \item Die TN werden verschiedene Märchen besprechen: fördert Toleranz \& kulturelle Sensibilität
    \item Die Fähigkeiten, kreativ zu sein, spontan zu sprechen und zu schreiben, werden geübt
    \item Mannigfaltige Aspekte der Kultur (Geschichte, Politik, Literatur usw.) werden auf einfachem Niveau mit Sprache verbunden
    \item Da es sich um Märchen handelt, werden die TN mehr Präteritum begegnen, d.h. sich oberflächlich damit bekannt machen
\end{itemize}

\large{\textbf{Aktivitäten und Aufgaben}}
\begin{enumerate}

    \item Fängt mit einem kurzen Video an: 
        \\\textit{KL sagt, dass die TN nicht unbedingt das ganze Video verstehen müssen. Sie müssen lediglich versuchen, so weit wie möglich beizubleiben.}
        \begin{itemize}
            \item \ulhref{https://youtu.be/JvDe3JtK6wI}{Lucia und die Bremer Stadtmusikanten} \par
            \textit{ggf.}
            \item \ulhref{https://youtu.be/eQvFp4ZD4Hc?list=RDeQvFp4ZD4Hc}{Was dast Du für große Ohren (Rotkäppchen)}
        \end{itemize}

    \item KL nennt kurz wie düster alte deutsche Märchen sind. z.B. „Suppenkaspar“ oder „Die Geschichte vom Bösen Friedrich“; „Aschenputtel“ ist auch dazu geeignet, weil sich die Brüder Grimm und Disney Versionen sehr stark unterscheiden

    \item Die TN bekommen eine Liste mit nützlichen Vokabeln und Redewendungen zum Thema „Märchen“ als auch eine ausgedrückte Version des Märchens   

    \item KL erwähnt kurz was Präteritum ist und warum es oft in Märchen auftaucht. Es soll auf die Verben „sein“ und „haben“ fokussiert werden. Eine Eselbrücke für die TN: Präteritum Verben haben oft „-te“ am Ende

    \item \textbf{PL:} KL bespricht kurz „Die Bremer Stadtmusikanten“ ggf. „Rotkäppchen“ anhand mancher Auzüge mit dem Beamer. KL soll erwähnen, dass die TN nicht alles wortwörtlich verstehen müssen, sondern nur versuchen zu im Groben zu folgen. Es wird hier auf Zeitformen fokussiert
        \begin{itemize}
            \item KL fasst den Inhalt des Märchens dadurch zusammen, die klassische Märchenstruktur (Anfang, Hauptteil, Ende) zu erwähnen. Dies ist für die Aufgabe später nützlich, da die TN selber einiges zum Märchen erfinden werden
        \end{itemize}

    \item \textbf{PL:} KL fragt die TN nach einigen Verben in dem Auszug im Präteritum und wie man es in Perfekt konjugiert. Falls TN Schwierigkeiten dabei haben: KL empfiehlt ihr, zuerst das Verb im Präsens und dann im Perfekt zu konjugieren

    \item \textbf{GA:} Die TN bilden dreiköpfige Gruppen. Jede Gruppe muss ein neues Ende für das im Unterricht behandeltes Märchen erfinden/das Märchen weiter fortsetzen nach dem Ende. Insgesamt $\pm$10 Sätze (im Präsens oder Perfekt)
    \begin{itemize}
        \item Jede Gruppe bekommt 4 Karten aus den jeweiligen Kategorien: Personen (wie z.B. der Bäcker, der Stiefmutter, der Fleischer usw.), Dinge (der Apfel, der Ring, der Schlüssel), Orte (der Fluss, die Heide, der Hügel) und Verben (entkommen, fressen, retten)
        \item Sie müssen jedes Wort bei der Aufgabe verwenden
        \item Sie können auch versuchen Wörter vom Handout beizufügen
    \end{itemize}

        \begin{itemize}
            \item Die TN können z.B. das Märchen heutzutage weitergehen lassen: Rotkäppchen trägt rote Jeans
        \end{itemize}

    \item \textbf{PL:} Einige Gruppen lesen vor. Fragen andere Gruppen darüber aus (fördert spontan sprechen)

\end{enumerate}

\large{\textbf{Im Fall von übrigen Zeit}}
\begin{itemize}
    
    \item \textbf{PA:} Die TN müssen kurz zur Zweit einander von ihrem Lieblingsmärchen erzählen. Möglichst auf Deutsch. Einige TN könnten die Erzählung ihres/ihrer Partners•in wiedergeben

\end{itemize}

\end{document}
